\scanpage{101}
\mspage{49}
% Main text, approximately lines 1--2.
In summary, we have a \deleted{\illegible{1}}\inserted{canonical} functor
% Approximately lines 3--10; the source of the vertical functor.
Isotopic category of \((X,S,\Gamma)\), \(X\) compact surface \deleted{\illegible{1}},
\(T\) oriented, \(S\) discrete subset, \(\Gamma\) operating by \unsure{\(\chi\)-automorphisms}
(\inserted{\(\sigma\in\Gamma\)}\unsure{preserving|preserves} the orientation if \(\chi(\sigma)=+1\), reversing it
if \(\chi(g)=-1\)), \(\Gamma\) transitive on \(\Pi_0(X)\),
\deleted{\illegible{1}} and if \((X_0,S_0)\) is a connected component
of \((X,S)\), we want that \deleted{\((X_0,S_0)\) \illegible{2}}\inserted{if \(X_0\simeq S^2\), we}
have \(\operatorname{card} S_{\unsure{0}}\geq 3\) \deleted{\illegible{3}}.
% Under "ait card S ≥ 3" a smaller line of text is struck through by one long stroke;
% "ait" itself is neither underlined nor struck.
\deleted{\illegible{1} \unsure{isomorphism}}
\margin{arrows: isotopy \unsure{classes} of homeomorphisms respecting all the structures.}

\diagram{\begin{tikzcd}[row sep=2.5em]
{} \arrow[d] \\
{}
\end{tikzcd}}

% Approximately lines 11--18.
Category of outer groups with \inserted{\(T\)-}loops \(\pi\),
\inserted{\(\pi\) noncommutative (i.e. \(\pi\neq 0,\mathbb Z\));}
on which a subgroup \(\Sigma\subset\Gamma\) of \(\Gamma\) \inserted{operates}
(this description being
equivalent: a description in terms
of systems of groupoids
\[
\Pi_{\unsure{D^*},\Gamma}\longrightarrow\Pi_{U,\Gamma}
\longrightarrow\Pi_{e,\Gamma}\quad\text{and }\kappa\ldots,
\]
more conceptual in certain contexts\ldots).
% Long oblique note in the left margin beside the target category.
\margin{Beware \unsure{of the} case \(S=\varnothing\) in the definition of the structures: \unsure{the} \illegible{1} (datum of
\(H^2(\pi,\mathbb Z)\overset{\kappa}{\simeq}T\)); we must of course take into account \unsure{during} the action of \(\Sigma\) on
\(H^2(\pi,\mathbb Z)\) \deleted{\unsure{and on \(T\)}}
(via its outer action on \(\pi\)) and on \(T\) (via \(\chi\)) and \unsure{require the} compatibility.}

% Approximately lines 19--25; the sentence continues on scan 102.
\deleted{\unsure{For definition}} I presume that the proof
\unsure{of the} fact that this functor is fully faithful
presents no essential difficulties,
% Oblique note in the lower-left margin, tied by a double bar and an arrow to the three preceding lines.
\margin{It would really be worth doing it very carefully\ldots}
using what is known for
\(\Gamma=1\). But the fact that, \unsure{for|\deleted{for}} \(\Gamma\) a given finite
group, \deleted{\unsure{that it is an isomorphism}}\inserted{it is \unsure{essential\textsuperscript{ly}} surjective}
% Lower-left handwritten pagination; attribution of this numbering is uncertain.
\margin{100}

\scanpage{102}
% Approximately lines 1--7; continuation from scan 101.
is an open problem on which
people are stumped. Of course we may assume \(\Sigma=\Gamma\),
\inserted{and \(\Gamma\subset\) extended Teichmüller}
and the question is whether every finite subgroup
of the extended Teichmüller is
realized as a group operating on \((X_0,S_0)\) \illegible{1}
in an essentially unique way. More
precisely, if \deleted{\illegible{1}} \(A=\) group of \deleted{\unsure{diffe}}homeomorphisms
% Approximately lines 8--16.
(or diffeomorphisms, if \(X_0\) is
differentiable) of \(X_0\), \(A^0\) its connected compo\unsure{nent|of}
(\uline{N.B.} \(A^0\) is contractible, \deleted{\unsure{)}}\inserted{in}
the anabelian cases) hence \(A/A^0=\mathcal T_{g,\nu}\)
(Teichmüller group for genus \(g\) and \(\nu\)
holes), the question comes down to: \unsure{this} whether for every
homomorphism from a finite group \(\Sigma\) into \deleted{\illegible{1}} \(\mathcal T_{g,\nu}\),
(we may assume \(\Sigma\subset\mathcal T_{g,\nu}\)), it lifts
to a homomorphism into \(A\), and whether two such
liftings are conjugate by an \(a\in A^0\)
(isotopy \inserted{in the strict sense} of two liftings \unsure{\(\sim\)}).

% Approximately lines 17--21; continuation on scan 103.
Having arrived at such a simple \deleted{description}\inserted{\unsure{reinterpretation|interpretation}}
of the 1-category (``1-isotopic'') deduced
from the 2-category of systems
\((\Pi_{D^*,\Gamma}\longrightarrow\Pi_{U,\Gamma}\longrightarrow\Pi_{e,\Gamma},\kappa)\), in terms
% Lower-left additional handwritten pagination; hand/status uncertain.
\margin{101}

\scanpage{103}
\mspage{50}
% Approximately lines 1--7; continuation from scan 102.
% The light round stroke before the square bracket may be the insertion mark for π.
of outer groups with loops (\inserted{\(\pi\)}[equipped
with a set of outer homomorphisms \(\kappa_i:T\longrightarrow\pi\), and
failing that with a \(\kappa:T\longrightarrow H^2(\pi,\mathbb Z)\))
(to do it properly, one should write
\(T^{(\otimes-1)}\simeq H^2(\pi,\mathbb Z)\), but here we have a canonical
isomorphism \(T^{\otimes-1}\simeq T\) i.e. \(T^{\otimes 2}\simeq\mathbb Z\ldots\))],
% Approximately lines 8--14.
the question arises \unsure{of how to recover},
(up to an equivalence defined up to unique isomorphism)
this diagram, in terms of \(\kappa\); everything
comes down to the \deleted{definition} description of the
categories \(\Pi_{D^*,\unsure{\Gamma}}\) and \(\Pi_{U,\Gamma}\), and the two
functors \(\Pi_{D^*,\Gamma}\longrightarrow\Pi_{U,\Gamma}\), \(\Pi_{U,\Gamma}\longrightarrow\Pi_{e,\Gamma}\);
\inserted{or what amounts to the \(=\), of the corresponding (multi-Galois) topoi and morphisms of topoi.}

% Approximately line 15.
a) \uline{Description} of \(\Pi_{U,\Gamma}\) and of \(\Pi_{U,\Gamma}\longrightarrow\Pi_{e,\Gamma}\).

% Approximately lines 16--22; continuation on scan 104.
Let more generally \(\pi\) be an \inserted{outer} group, \deleted{\unsure{such that}} \unsure{whose}
\uline{center} \uline{is trivial} (this corresponds to the
anabelian hypothesis!), let us show how we
associate with it a classifying topos
\(B_\pi\), which (as a category of sheaves) will be
% "(pour Iso)" is written under "faisceaux) sera" and tied by a curve to the place after "fonctorielle".
\(\operatorname{Ens}(\pi)\)) in a ``functorial'' way \inserted{(for \unsure{Isomorphisms|isomorphisms})}. Everything
comes down to seeing how, to a \uline{conjugacy}
% Lower-left additional handwritten pagination; hand/status uncertain.
\margin{102}

\scanpage{104}
% Approximately lines 1--7; continuation from scan 103.
\uline{class} \inserted{\(\Theta\)} of isomorphisms
\[
\pi'\xrightarrow{u}\pi,
\]
we associate an ``inverse image'' functor
\(\operatorname{Ens}(\pi)\longrightarrow\operatorname{Ens}(\pi')\), defined up to unique
isomorphism. Let us consider for every \(u\) in the class
\(\Theta\), the functor ``\(u\)-restriction of operations''
\[
\operatorname{Ens}(\pi)\xrightarrow{u^*}\operatorname{Ens}(\pi')
\]
% Approximately lines 8--11, including interlinear corrections.
% "définit," stands above the struck "est"; "avec ... ," stands above the line and is tied by a curve to the place after "donc,".
[which \deleted{\unsure{is}}\inserted{defines,} therefore, \inserted{with \(\unsure{u_*=}(u^{-1})^*\),}
an \deleted{\unsure{isomorphism}}\inserted{equivalence} of the Topoi
\[
B_{\pi'}\xrightarrow{\substack{u_\bullet\\ \sim}} B_\pi
\]
]
We shall, between these \deleted{functors}\inserted{equivalences}, \inserted{for \(u\in\Theta\),} define a
transitive system of isomorphisms (which therefore lets us
identify them with each other!).

% Approximately lines 12--20.
\unsure{Let|Let} \(u,u'\in\Theta\), whence \(u^*,u'^*\); we have by hypothesis
\deleted{definition} a \(g\in\pi\) such that \(u'=\operatorname{int}(g)\circ u\),
furthermore \(g\) is determined modulo an
element of \(\operatorname{Centr}_\pi u(\pi')=\operatorname{Centr}(\pi)=1\);
so here \(g\) is unique. But \(g\)
can be used to define a functorial
isomorphism \deleted{between \(u^*\) and \(u'^*\)}\inserted{between}
\[
i_{u',u}:u^*\xrightarrow{\sim}u'^*
\]
by taking, for \(E\in\operatorname{Ob}\operatorname{Ens}(\pi)\),
\[
i_{u',u}(E):
\underset{\substack{\unsure{\Vert}\\E}}{u^*(E)}
\longrightarrow
\underset{\substack{\unsure{\Vert}\\E}}{u'^*(E)},
\qquad i_{u',u}(E)=g_E.
\]
% Lower-left additional handwritten pagination; hand/status uncertain.
\margin{103}

\scanpage{105}
\mspage{51}
% Approximately lines 1--6.
\uline{N.B.} The argument works for every conjugacy
class of homomorphisms \inserted{of groups} \(\pi'\longrightarrow\pi\) (i.e. every outer
homomorphism) whose centralizer in \(\pi\) is
reduced to 1 (which for an epimorphism
\unsure{reduces to the} hypothesis \(\operatorname{Centr}(\pi)=1\)).
% Oblique note in the left margin. The word read "anabéliens" carries a double rule through
% or under it; deletion or emphasis is not settled.
\margin{works for \unsure{anabelian loops|anabelian groups with loops}, and the homomorphisms of such
whose image is of finite index\ldots{} (for the centralizer in such a group \(\pi\)
of a subgroup of finite \unsure{index} is still \unsure{reduced to e|reduced to 1}).}

% Approximately lines 7--14.
If now a group \(\Sigma\) operates
on the outer group \inserted{\(\pi\)}, then by the
preceding result we can say that it
also operates on the topos \(B_\pi\), \inserted{whence a topos \(B_{\pi,\Gamma}\),} (we can
also say that \deleted{\illegible{1}} Since
\(\operatorname{Centr}(\pi)=1\)) the operation of \deleted{\(\unsure{\pi}\)}\(\Sigma\) on \(\pi\)
defines an extension \(\mathcal E\) of \(\Sigma\) by \(\pi\),
% "On récupère" stands above the struck word; a curve runs from it over "bien" to the end of the line.
whence a topos \(B_{\pi,\Gamma}=B_{\mathcal E}\), \deleted{\unsure{and}} \inserted{We recover} \unsure{of}
course also \(B_{\pi,\Gamma}\longrightarrow B_{\Gamma}\). To reassure
% Approximately lines 15--22; continuation on scan 106.
ourselves, we should make sure that
if we have a homomorphism of outer groups
\(\pi'\xrightarrow{\Theta}\pi\) commuting with the action of a
\(\Sigma\), with \(\operatorname{Centr}_\pi(\Theta)=(1)\), then there exists
a \(\Theta_{\Gamma}:B_{\pi',\Gamma}\longrightarrow B_{\pi,\Gamma}\) defined up to canonical
isomorphism. Now every \(u\in\Theta\) defines a
homomorphism of extensions \inserted{\(u_{\Gamma}\) (of \(\Sigma\) by \(\pi'\) resp. \(\pi\))}
\(\mathcal E'\longrightarrow\mathcal E\) \unsure{above|below}
\(u\), and if we pass from \(u\) to \(u'=\operatorname{int}(g)\circ u\), we will have
% Lower-left additional handwritten pagination; hand/status uncertain.
\margin{104}